% ─────────────────────────────────────────────────────────────────────────────
%  SECTION 3 — THE BAILOUT PROTOCOL
%  Formal specification: state machine, trigger conditions, escalation algorithm,
%  bypass mechanism, and timeout handling.
% ─────────────────────────────────────────────────────────────────────────────

\section{The Bailout Protocol}\label{sec:bailout}

The Bailout Protocol is the mandatory exception-escalation mechanism embedded in
every node of a \bxthree{} system. It is not an error-handling routine selected
at runtime, nor an optional escalation pathway available to a system that chooses
to invoke it. It is an \emph{architectural compulsion}: any condition a node
cannot resolve with certainty must propagate to a Human Accountability Anchor,
and no machine actor in the propagation chain may absorb, resolve, or suppress
that propagation. This section provides the complete formal specification:
the deterministic state machine governing exception propagation, the precise
conditions governing each transition, the bail-out trigger condition, the
escalation algorithm, the bypass mechanism, and the timeout hierarchy that
guarantees liveness under all failure conditions.

% ─────────────────────────────────────────────────────────────────────────────
%  3.1  State Machine Formal Definition
% ─────────────────────────────────────────────────────────────────────────────

\subsection{State Machine Definition}\label{sec:bailout-state-machine}

The Bailout Protocol is formally modeled as a deterministic finite state machine
embedded in the \fact{} layer of every \bxthree{} node:

\[
\mathcal{S}_{\text{bailout}} \;=\; \bigl( Q,\; \Sigma,\; \delta,\; q_0,\; F \bigr)
\]

\begin{itemize}\itemsep=0pt
\item[$Q$] $\{\;\text{IDLE},\; \text{BOUNDED},\; \text{BAIL\_PENDING},\;
  \text{ESCALATING},\; \text{HUMAN\_ACKNOWLEDGED},\; \text{RESOLVED}\; \}$ is the
  finite set of protocol states.
\item[$\Sigma$] The input alphabet defined in Table~\ref{tab:bailout-events}.
\item[$\delta : Q \times \Sigma \to Q$] The deterministic transition function
  (Table~\ref{tab:bailout-transition}). For every $(q, \sigma) \in Q \times
  \Sigma$, at most one $q' \in Q$ satisfies $\delta(q, \sigma) = q'$.
  If no transition is defined, the machine stalls and a protocol violation is
  recorded in the Forensic Ledger.
\item[$q_0$] $\text{IDLE}$ is the designated initial state, entered at node
  startup and after every successful directive resolution.
\item[$F$] $\{\;\text{RESOLVED}\; \}$ is the set of terminal accepting states.
\end{itemize}

A node's state $q \in Q$ is stored in the node's \fact{} layer worksheet and is
atomically observable by its parent node and by the Bailout Monitor. State updates
are atomic: no intermediate or partially-executing states are externally
observable. The machine is \emph{deterministic} by construction; the Bailout
Monitor's priority function $\pi$ totally orders simultaneous trigger conditions
and selects the highest-priority enabled transition before committing any state
change to the Forensic Ledger.

\bigskip

\noindent\textbf{State definitions:}

\begin{trivlist}\itemindent=2em
\item[\textbf{IDLE}] The node's Bounds Engine is operating within its provisioned
  capability. No exceptional condition exists. Normal sensing, modelling, and
  proposition cycles are active.

\item[\textbf{BOUNDED}] The node has detected a condition that \emph{may} require
  escalation, but the Bounds Engine is still evaluating whether the condition is
  resolvable locally or requires human intervention. The condition is known and
  contained. No propagation has occurred.

\item[\textbf{BAIL\_PENDING}] The node has determined the condition cannot be
  resolved locally. The trigger package $\mathcal{T}$ has been assembled and is
  in the process of being propagated to the parent node. Delivery is asynchronous
  and confirmed.

\item[\textbf{ESCALATING}] The trigger package has been delivered to a machine
  actor in the propagation chain. That actor is evaluating whether it can resolve
  the exception. If it cannot, it propagates further upward. This state persists
  across zero or more machine-actor hops.

\item[\textbf{HUMAN\_ACKNOWLEDGED}] The trigger package has reached a Human
  Accountability Anchor. The human operator has received the alert and is actively
  reviewing the exception. The system awaits a directive from the human.

\item[\textbf{RESOLVED}] The human has issued a directive and the directive has
  propagated back to the originating node. The actuators have executed or will
  execute the directive. The trigger is closed and the complete lifecycle is
  sealed in the Forensic Ledger.
\end{trivlist}

% ─────────────────────────────────────────────────────────────────────────────
%  3.2  Transition Conditions and Transition Table
% ─────────────────────────────────────────────────────────────────────────────

\subsection{Transition Conditions}\label{sec:bailout-transitions}

Let $e \in \Sigma$ denote the event driving a transition. Table
\ref{tab:bailout-transition} defines $\delta(q, e)$. Entries marked ``—''
are undefined transitions: the machine stalls and a protocol violation is logged.

\begin{table}[htbp]
\centering
\caption[Bailout Protocol state transition table.]
{Bailout Protocol state transition table. $\checkmark$ denotes a defined
transition; ``---'' denotes undefined (stall-and-violate).}
\label{tab:bailout-transition}
\bigskip
\begin{tabular}{@{}lcccccc@{}}
\toprule
 & \multicolumn{6}{c}{Input event $e$} \\
\cmidrule(l){2-7}
Current state & $E_{\text{anomaly}}$ & $E_{\text{local\_ok}}$ &
$E_{\text{bail}}$ & $E_{\text{prop}}$ & $E_{\text{ack}}$ &
$E_{\text{directive}}$ \\
\midrule
IDLE                  & BOUNDED           & ---            & ---           & ---            & ---            & --- \\
BOUNDED               & ---               & IDLE           & BAIL\_PENDING & ---            & ---            & --- \\
BAIL\_PENDING         & ---               & ---            & ---           & ESCALATING     & ---            & --- \\
ESCALATING            & ---               & ---            & ---           & ESCALATING     & HUMAN\_ACK     & RESOLVED \\
HUMAN\_ACKNOWLEDGED   & ---               & ---            & ---           & ---            & IDLE           & RESOLVED \\
RESOLVED              & ---               & ---            & ---           & ---            & ---            & ---
\end{tabular}
\end{table}

\bigskip
\noindent\textbf{Event definitions ($E \subset \Sigma$):}

\begin{trivlist}\itemindent=2em
\item[$E_{\text{anomaly}}$] An input or sensor reading falls outside the
  node's nominal operating range $R(n)$, or an explicit \texttt{bailout\_signal}
  is emitted by the Bounds Engine.

\item[$E_{\text{local\_ok}}$] The Bounds Engine evaluates the condition and
  determines it can be resolved within the current capability scope without
  human intervention.

\item[$E_{\text{bail}}$] The Bounds Engine confidence score $\phi$ for local
  resolution falls below the provisioned threshold $\theta_{\text{bail}}$:
  $\phi < \theta_{\text{bail}}$.

\item[$E_{\text{prop}}$] The trigger package has been successfully delivered
  to the next node in the accountability chain (parent or human anchor).

\item[$E_{\text{ack}}$] The human operator has acknowledged receipt of the
  trigger package and is actively reviewing it.

\item[$E_{\text{directive}}$] The human operator has issued a binding directive
  resolving the exception. The directive is authoritative and supersedes any
  machine-generated proposal.
\end{trivlist}

\bigskip
\noindent\textbf{Guard conditions:} Each transition in Table
\ref{tab:bailout-transition} is governed by a guard condition evaluated by the
\fact{} layer pre-commit hook before the transition is committed:

\begin{itemize}\itemsep=0pt
\item $\text{IDLE} \xrightarrow{E_{\text{anomaly}}} \text{BOUNDED}$:
  $\text{TRIGGER}(n, x)$ fires for anomaly $x$ (defined in
  Section~\ref{sec:bailout-trigger}).

\item $\text{BOUNDED} \xrightarrow{E_{\text{bail}}} \text{BAIL\_PENDING}$:
  $\phi(x) \geq \theta_{\text{bail}}$ for the strongest unresolved condition
  $x$ in the BOUNDED state.

\item $\text{BAIL\_PENDING} \xrightarrow{E_{\text{prop}}} \text{ESCALATING}$:
  Delivery confirmation received from the parent node; propagation chain entry
  written to Forensic Ledger.

\item $\text{ESCALATING} \xrightarrow{E_{\text{ack}}} \text{HUMAN\_ACKNOWLEDGED}$:
  $\text{isHuman}(\text{current}) = \text{true}$ and acknowledgment signal
  received at the human interface.

\item $\text{ESCALATING} \xrightarrow{E_{\text{directive}}} \text{RESOLVED}$:
  Directive received at originating node; actuator command issued or queued.

\item $\text{HUMAN\_ACKNOWLEDGED} \xrightarrow{E_{\text{directive}}}
  \text{RESOLVED}$: Directive received; actuator command issued; trigger package
  closed.
\end{itemize}

% ─────────────────────────────────────────────────────────────────────────────
%  3.3  Bail-Out Trigger Condition
% ─────────────────────────────────────────────────────────────────────────────

\subsection{Bail-Out Trigger Condition}\label{sec:bailout-trigger}

The Bailout Protocol is activated by a single, general trigger condition:

\begin{equation}
\text{TRIGGER}(n, x) \;:=\; \bigl( x \notin R(n) \bigr) \;\lor\;
  \bigl( \phi(x) < \theta_{\text{bail}} \bigr) \;\lor\;
  \texttt{bailout\_signal}(x) \tag{TC}
\end{equation}

Where:

\begin{itemize}\itemsep=0pt
\item $n$ is the originating node.
\item $x$ is the condition (input anomaly, state conflict, or safety
  envelope prediction) under evaluation.
\item $R(n)$ is the operating range provisioned for $n$ by the \purpose{}
  layer at node initialization.
\item $\phi(x)$ is the Bounds Engine's confidence score for resolving $x$
  locally, computed from the evidence available to the \bounds{} layer.
\item $\theta_{\text{bail}}$ is the provisioned confidence threshold below
  which local resolution is rejected. It is set at node provisioning and is
  immutable at runtime; no machine actor may adjust $\theta_{\text{bail}}$ to
  suppress escalation.
\item $\texttt{bailout\_signal}(x)$ is an explicit signal emitted by the
  Bounds Engine when it detects a logical contradiction or unresolvable
  state within $R(n)$ that does not reduce to a low-confidence prediction.
\end{itemize}

\bigskip
\noindent\textbf{Three canonical trigger subclasses:} Condition (TC) covers
three distinct trigger modes:

\begin{enumerate}\itemsep=0pt
\item \textbf{Capability Boundary}: $x \notin R(n)$ because $x$ falls outside
  the provisioned scope of $n$. Example: a soil moisture node receives a
  command to adjust fertilizer composition — a function outside its
  provisioned operating range.

\item \textbf{Safety Envelope Violation (predicted)}: The Bounds Engine's
  forward simulation predicts that executing a proposed action would violate
  a safety constraint $s \in S(n)$, where $S(n)$ is the safety envelope
  provisioned by the \purpose{} layer.

\item \textbf{Accountability Boundary Violation}: A proposed action $a$
  requires authority that no machine actor in the propagation chain possesses
  (e.g., signing a financial commitment, overriding a regulatory parameter,
  accepting liability). This trigger mode fires when $a$ enters the escalation
  path and no machine actor can claim terminal resolution authority.
\end{enumerate}

\bigskip
\noindent\textbf{The no-discretion property:} The \bounds{} layer does not
exercise discretion over the trigger condition. Specifically:

\begin{enumerate}
\item The evaluation of $\text{TRIGGER}(n, x)$ is a read-only query against
  the operating range definition $R(n)$, not a judgment call made by the
  reasoning system.
\item The result of the query directly drives the state machine transition
  without an intervening decision step. There is no machine-actor gate
  between trigger evaluation and protocol activation.
\item The confidence threshold $\theta_{\text{bail}}$ is provisioned at
  node initialization and is immutable at runtime. No machine actor may
  adjust $\theta_{\text{bail}}$ to suppress escalation.
\end{enumerate}

These constraints ensure that the trigger condition is architecturally
enforced rather than operationally optional.

% ─────────────────────────────────────────────────────────────────────────────
%  3.4  Escalation Algorithm
% ─────────────────────────────────────────────────────────────────────────────

\subsection{Escalation Algorithm}\label{sec:bailout-escalation}

Upon entering the ESCALATING state, the node executes the following escalation
algorithm to route the exception to the Human Accountability Anchor. The
algorithm is embedded identically in every node's Worksheet; it is
deterministic and cannot be overridden by any machine actor.

\bigskip
\begin{algorithm}[H]
\caption{Bailout Protocol Escalation Algorithm (executed at every node)}
\label{alg:bailout-escalation}
\begin{enumerate}
\item \textbf{Initialize.} Set $\text{current} \leftarrow n$,
  $\text{propagation\_chain} \leftarrow []$, $r \leftarrow 0$,
  $\text{hop\_count} \leftarrow 0$.

\item \textbf{Check terminus.} If $\text{isHuman}(\text{current}) = \text{true}$,
  go to Step~\ref{step:deliver}.

\item \textbf{Record propagation step.} Append
  $\{\text{node}: \text{current.id},\; \text{action}: \texttt{ESCALATING},\;
  t_{\text{now}},\; \text{hop}: \text{hop\_count}\}$ to
  $\text{propagation\_chain}$.

\item \textbf{Evaluate at machine actor.} Set
  $r \leftarrow \text{current.BoundsEngine.evaluate}(\mathcal{T})$.
  \begin{enumerate}
  \item If $r.\text{can\_resolve} = \text{true}$ and
    $r.\text{within\_SafetyEnvelope} = \text{true}$:
    \begin{enumerate}
    \item $\mathcal{T}.\text{resolution} \leftarrow r$
    \item $\mathcal{T}.\text{status} \leftarrow \texttt{RESOLVED\_LOCAL}$
    \item Propagate resolution back down the chain to $n$; \textbf{return}
    \end{enumerate}
  \item Else: continue to Step~\ref{step:parent}.
  \end{enumerate}

\item \label{step:parent}\textbf{Move to parent.} Set $\text{current}
  \leftarrow \alpha(\text{current})$ (the immutable parent pointer).
  Increment $\text{hop\_count}$.

\item \textbf{Hard hop limit.} If $\text{hop\_count} > \text{HOP\_LIMIT}$
  (default: 127), transition the originating node to
  $\texttt{QUIESCENT}$ and alert the Human Root via fallback channels;
  record $\texttt{ESCALATION\_LOOP\_PREVENTED}$ in the Forensic Ledger;
  \textbf{return}.

\item \textbf{Retry with backoff.} Set $\text{delay} \leftarrow
  \text{Backoff}(\text{hop\_count}, T_{\text{prop}})$; wait \text{delay};
  attempt delivery to $\text{current}$. If delivery fails, go to
  Step~\ref{step:parent}.

\item \label{step:deliver}\textbf{Deliver to human.} Set
  $\mathcal{T}.\text{status} \leftarrow \texttt{HUMAN\_ACKNOWLEDGED}$.
  Deliver $\mathcal{T}$ to the human interface (Root Dashboard alert,
  SMS, or fallback). Await directive $d$ from the human operator.

\item \textbf{Process human response.} If $d$ received before
  $T_{\text{human}}$:
  \begin{enumerate}
  \item $\mathcal{T}.\text{directive} \leftarrow d$
  \item $\mathcal{T}.\text{status} \leftarrow \texttt{RESOLVED}$
  \item Propagate $d$ back down the propagation chain to $n$; \textbf{return}
  \end{enumerate}
  If $T_{\text{human}}$ expires without response:
  \begin{enumerate}
  \item Transition originating node to $\texttt{QUIESCENT}$
  \item Alert Human Root via all available fallback pathways
  \item $\mathcal{T}.\text{status} \leftarrow \texttt{RESOLVED\_TIMEOUT}$
  \item \textbf{return}
  \end{enumerate}
\end{enumerate}
\end{algorithm}

\bigskip
\noindent\textbf{Backoff function:} The propagation backoff between parent
hops is computed by the Bailout Monitor and follows an exponential curve:

\begin{equation}
\text{Backoff}(k, T_{\text{prop}}) \;=\; \min\Bigl(
  T_{\text{prop}} \cdot 2^{k},\;
  T_{\text{cap}} \Bigr) \tag{eq:backoff}
\end{equation}

Where $k$ is the current hop count and $T_{\text{prop}}$ is the base
propagation interval. This backoff is defined in the protocol specification
and cannot be altered at runtime by any machine actor, preventing
denial-of-escalation attacks that would otherwise use artificially short
backoff intervals to overwhelm the human operator.

% ─────────────────────────────────────────────────────────────────────────────
%  3.5  Bypass Mechanism
% ─────────────────────────────────────────────────────────────────────────────

\subsection{Bypass Mechanism: Why the Protocol Bypasses All Machine Actors}
\label{sec:bailout-bypass}

The Bailout Protocol's bypass of all intermediate machine actors is not an
arbitrary design choice — it is a direct consequence of the authority
structure of the \bxthree{} framework. A Machine Accountability Anchor holds:
proposal authority (can propose state transitions), simulation authority (can
model outcomes in sandbox), zero execution authority (cannot act on the
physical world), and zero legal authority (cannot sign, commit, or authorize
on behalf of the organisation).

When a trigger fires, the requested resolution typically requires at least one
of: overriding a Safety Envelope parameter, authorising a financial
commitment, making a regulatory determination, or accepting liability for an
action's outcome. None of these are within a machine actor's authority
class. A machine actor that absorbed a trigger rather than propagating it
would be claiming authority it does not possess — structurally
indistinguishable from a compromised or rogue node.

\bigskip
\noindent\textbf{Machine Actor Exclusion (MAE) rule:}
\begin{invariant}[MAE]
  No machine actor may serve as a terminal node in the Bailout Protocol
  escalation chain. The only admissible terminal nodes are Human
  Accountability Anchors. This property is enforced by the \fact{} layer
  gate: a machine actor that receives a trigger package must either
  propagate it upward or deadlock.
\end{invariant}

The MAE rule is enforced by the transition function $\delta$:
the ESCALATING $\xrightarrow{E_{\text{prop}}}$ ESCALATING transition
is the only defined transition available at a machine actor; no transition
exists from ESCALATING to RESOLVED through a machine actor as terminus.
A machine actor that cannot propagate (e.g., its parent pointer is null)
must enter a deadlock state and alert the Human Root via fallback channels.

\bigskip
\noindent\textbf{Failure modes eliminated by the bypass:} The bypass
mechanism structurally eliminates three failure modes that are endemic to
any architecture in which machine actors can serve as terminal escalation
nodes:

\begin{itemize}\itemsep=0pt
\item \textbf{Intercept}: A machine actor suppresses the trigger and handles
  it autonomously without human involvement. Bypass makes this impossible
  — the machine actor has no defined transition from ESCALATING to
  RESOLVED.

\item \textbf{Redirect}: A machine actor routes the trigger to a different
  machine actor rather than the designated human anchor. Bypass makes
  this structurally impossible — the propagation path is determined by the
  immutable parent pointer chain $\alpha(\cdot)$, not by any routing table
  under machine-actor control.

\item \textbf{Delay}: A machine actor defers propagation pending ``further
  analysis.'' Bypass makes this architecturally infeasible — the \fact{}
  layer pre-commit hook enforces $T_{\text{bounds}}$ on the BOUNDED state,
  and the propagation backoff curve is defined in the protocol specification
  and enforced by the Bailout Monitor, not by the machine actor.
\end{itemize}

\bigskip
\noindent\textbf{Edge case — Human Accountability Anchor unreachable:} If
the human is unreachable (off-grid, medical emergency, disconnected), the
system does not continue operating autonomously. It halts gracefully in
QUIESCENT state, retaining all Forensic Ledger entries and awaiting
explicit human reset. This is a deliberate design decision: the system
prioritises never acting without human authorisation over operational
continuity.

% ─────────────────────────────────────────────────────────────────────────────
%  3.6  Timeout Handling
% ─────────────────────────────────────────────────────────────────────────────

\subsection{Timeout Handling}\label{sec:bailout-timeouts}

The Bailout Protocol defines four independent timeouts, each governing a
distinct phase of the protocol and enforced by a different component:

\bigskip
\begin{table}[htbp]
\centering
\caption{Bailout Protocol timeout hierarchy}
\label{tab:bailout-timeouts}
\bigskip
\begin{tabular}{@{}lp{5cm}p{5cm}@{}}
\toprule
\textbf{Timeout} & \textbf{Governs} & \textbf{Enforced by} \\
\midrule
$T_{\text{bounds}}$ & BOUNDED state duration; time allowed for the
  Bounds Engine to evaluate and determine local resolvability &
  \fact{} layer pre-commit hook \\
$T_{\text{prop}}$ & Propagation interval between parent hops;
  time between delivery attempt and confirmation &
  Bailout Monitor, enforced independently of \bounds{} \\
$T_{\text{cap}}$ & Pathway capability window; time after which a
  failed pathway is marked \texttt{DEGRADED} and an alternate is selected &
  Pathway Health Registry \\
$T_{\text{human}}$ & Human acknowledgment window; time the system awaits
  a directive from the human after \texttt{HUMAN\_ACKNOWLEDGED} &
  Human interface layer; hard halt on expiry \\
\midrule
\end{tabular}
\end{table}

\bigskip
\noindent\textbf{Bounds timeout $T_{\text{bounds}}$:} The BOUNDED state
has a maximum duration $T_{\text{bounds}}$, enforced by the \fact{} layer
pre-commit hook. If the Bounds Engine does not emit either
$E_{\text{local\_ok}}$ or $E_{\text{bail}}$ before $T_{\text{bounds}}$
expires, the machine transitions automatically to BAIL\_PENDING via
$E_{\text{timeout}}$. This prevents the Bounds Engine from delaying
escalation indefinitely while operating on an unresolved condition.
$T_{\text{bounds}}$ is a deployment parameter provisioned at node
initialisation and is immutable at runtime.

\bigskip
\noindent\textbf{Propagation timeout $T_{\text{prop}}$:} Each propagation
hop between parent nodes is governed by $T_{\text{prop}}$. If delivery to
the parent is not confirmed within this window, the Bailout Monitor
re-attempts with exponential backoff per Equation~(eq:backoff). After
$\text{HOP\_LIMIT}$ hops with no confirmation, the system enters
QUIESCENT and alerts the Human Root via fallback channels.

\bigskip
\noindent\textbf{Pathway health timeout $T_{\text{cap}}$:} The Pathway
Health Registry (PHR) tracks the availability and latency of each
provisioned communication pathway to the human anchor. Before each
propagation step, the PHR is queried and the lowest-latency functional
pathway is selected. Pathways that fail $K$ consecutive health checks
(default $K = 3$) are marked \texttt{DEGRADED} and excluded from
selection until a health check succeeds. A background health-ping thread
updates the PHR at intervals of $T_{\text{health}}$ (default: 30 seconds).

\bigskip
\noindent\textbf{Human acknowledgment timeout $T_{\text{human}}$:} After
the trigger package reaches the Human Accountability Anchor and the node
enters HUMAN\_ACKNOWLEDGED, the system awaits a directive for a maximum
of $T_{\text{human}}$. If no directive is received:

\begin{enumerate}
\item The system transitions to QUIESCENT — a strict subset of RESOLVED
  in which the node suspends all actuator commands and awaits explicit
  human reset.
\item The Human Root is alerted via all available fallback pathways.
\item The trigger is recorded as $\texttt{RESOLVED\_TIMEOUT}$ in the
  Forensic Ledger, ensuring the failure is attributed as a human
  non-response incident requiring post-mortem, not a silent suppression.
\end{enumerate}

\bigskip
\noindent\textbf{Worst-case liveness bound:} Together, the four timeouts
define a finite upper bound on wall-clock time from trigger firing (T1) to
human acknowledgment under any combination of parent unavailability and
pathway degradation:

\begin{equation}
T_{\max} \;=\; T_{\text{bounds}} \;+\;
N_{\max} \cdot \bigl( T_{\text{prop}} \cdot 2^{\text{HOP\_LIMIT}} \bigr)
\;+\; T_{\text{human}} \tag{eq:tmax}
\end{equation}

Where $N_{\max}$ is the maximum delegation depth (derived from the network
topology). A deployment that cannot guarantee human acknowledgment within
$T_{\max}$ must provision additional or higher-bandwidth pathways, or must
reduce $T_{\text{prop}}$ to bring the worst-case bound within the required
limit. This analysis is a required part of deployment certification for any
\bxthree{} system operating in a regulated or safety-critical environment.

\bigskip
\noindent\textbf{Timeout interaction with the bypass mechanism:} The bypass
mechanism ensures that timeouts cannot be weaponised by a machine actor to
delay or suppress escalation:

\begin{itemize}\itemsep=0pt
\item $T_{\text{bounds}}$ is enforced by the \fact{} layer pre-commit hook,
  not by the \bounds{} layer. A machine actor cannot artificially extend
  $T_{\text{bounds}}$ to prevent transition to ESCALATING.
\item $T_{\text{prop}}$ backoff is calculated by the Bailout Monitor, not
  by any machine actor. The backoff curve is defined in the protocol
  specification and cannot be altered at runtime.
\item If all pathways timeout at $T_{\text{cap}}$, the protocol resolves
  with $\texttt{PATHWAY\_EXHAUSTED}$ — a terminal RESOLVED state, not a
  suppression of escalation. The tag in the Forensic Ledger ensures the
  failure is recorded as an incident, not a silent suppression.
\end{itemize}

% ─────────────────────────────────────────────────────────────────────────────
%  3.7  Summary of Properties
% ─────────────────────────────────────────────────────────────────────────────

\subsection{Summary of Protocol Properties}\label{sec:bailout-summary}

The Bailout Protocol satisfies the following formal properties:

\begin{enumerate}\itemsep=0pt
\item $\mathcal{S}_{\text{bailout}}$ is a deterministic finite state machine
  with six states and a single terminal state (RESOLVED).

\item The trigger condition TC is individually necessary and jointly sufficient
  to cover all conditions requiring human escalation: capability boundary,
  safety envelope prediction, and accountability boundary violation.

\item The Machine Actor Exclusion rule (Invariant MAE) is enforced by the
  transition function $\delta$ and cannot be disabled or relaxed by any
  machine actor. The bypass mechanism eliminates the Intercept, Redirect,
  and Delay failure modes by architectural constraint.

\item The escalation algorithm (Algorithm~\ref{alg:bailout-escalation})
  terminates within $\text{distance}(n, h(n)) + 1$ propagation hops, or triggers
  QUIESCENT halt on pathway exhaustion.

\item Every state transition is immutably recorded in the Forensic Ledger
  with a SHA-256 digest chain for tamper-evident auditability.

\item The four-level timeout hierarchy ($T_{\text{bounds}}$,
  $T_{\text{prop}}$, $T_{\text{cap}}$, $T_{\text{human}}$) guarantees a
  finite worst-case liveness bound $T_{\max}$ for every protocol run.

\item The bypass mechanism guarantees that human reachability survives
  delegation depth, network partition, and recursive node spawning, by
  exploiting the immutable parent-pointer structure and human accountability
  anchor reference embedded in each node's worksheet.
\end{enumerate}

These properties collectively ensure that the Bailout Protocol delivers on
its architectural guarantee: \emph{any} condition a machine actor cannot
resolve with certainty \emph{must} propagate to a Human Accountability
Anchor, and no machine actor in the propagation chain may suppress,
absorb, or delay that propagation.
